% ============================================================
% PREAMBLE.TEX
% Распределённая оптимизация и машинное обучение
% ============================================================

% ------------------------------------------------------------
% МАТЕМАТИЧЕСКИЕ ПАКЕТЫ
% ------------------------------------------------------------

\usepackage{amsmath}
\usepackage{amssymb}
\usepackage{amsthm}
\usepackage{mathtools}
\usepackage{bm}
\usepackage{mathrsfs}

% Дополнительные математические символы и интегралы
\usepackage{esint}

% Удобная работа с условиями и командами
\usepackage{etoolbox}

% ------------------------------------------------------------
% ТИПОГРАФИКА
% ------------------------------------------------------------

\usepackage{microtype}
\usepackage{ragged2e}
\usepackage{enumitem}

% Красивое оформление заголовков
\usepackage{titlesec}
\usepackage{tocloft}

% Колонтитулы
\usepackage{fancyhdr}

% Улучшенные подписи
\usepackage{caption}
\usepackage{subcaption}

% ------------------------------------------------------------
% ТАБЛИЦЫ
% ------------------------------------------------------------

\usepackage{booktabs}
\usepackage{array}
\usepackage{tabularx}
\usepackage{longtable}
\usepackage{multirow}

% ------------------------------------------------------------
% РИСУНКИ И ГРАФЫ
% ------------------------------------------------------------

\usepackage{tikz}

\usetikzlibrary{
    arrows.meta,
    calc,
    positioning,
    fit,
    shapes.geometric,
    shapes.misc,
    decorations.pathreplacing,
    decorations.markings,
    patterns,
    backgrounds,
    graphs,
    quotes,
    matrix
}

\graphicspath{
    {figures/}
    {images/}
    {./}
}

% ------------------------------------------------------------
% АЛГОРИТМЫ И ПСЕВДОКОД
% ------------------------------------------------------------

\usepackage{float}
\usepackage{algorithm}
\usepackage{algpseudocode}

% ------------------------------------------------------------
% ЛИСТИНГИ ПРОГРАММНОГО КОДА
% ------------------------------------------------------------

\usepackage{listings}

% ------------------------------------------------------------
% ССЫЛКИ
% hyperref уже загружен в main.tex
% ------------------------------------------------------------

\usepackage{xurl}
\usepackage{bookmark}
\usepackage[nameinlink,noabbrev]{cleveref}

% ============================================================
% ЦВЕТА
% ============================================================

\definecolor{MaulRed}{HTML}{8A1010}
\definecolor{MaulDarkRed}{HTML}{5A0808}
\definecolor{TextBlack}{HTML}{1E1E1E}
\definecolor{DarkGray}{HTML}{444444}
\definecolor{MiddleGray}{HTML}{777777}
\definecolor{LightGray}{HTML}{F2F2F2}
\definecolor{CodeGray}{HTML}{F6F6F6}

% ============================================================
% ОБЩАЯ ТИПОГРАФИКА
% ============================================================

\setlength{\parindent}{1.25cm}
\setlength{\parskip}{0pt}

\setlength{\emergencystretch}{2em}

\clubpenalty=10000
\widowpenalty=10000
\displaywidowpenalty=10000

\setcounter{secnumdepth}{3}
\setcounter{tocdepth}{3}

% Ширина полей нумерации в оглавлении.
\setlength{\cftchapnumwidth}{2.8em}
\setlength{\cftsecnumwidth}{3.8em}
\setlength{\cftsubsecnumwidth}{4.8em}
\setlength{\cftsubsubsecnumwidth}{5.8em}

\allowdisplaybreaks[2]

% Немного увеличиваем расстояние между строками в формулах
\setlength{\jot}{5pt}

% Расстояния между плавающими объектами и текстом
\setlength{\textfloatsep}{18pt plus 2pt minus 4pt}
\setlength{\floatsep}{14pt plus 2pt minus 3pt}
\setlength{\intextsep}{14pt plus 2pt minus 3pt}

% ============================================================
% РУССКИЕ НАЗВАНИЯ
% ============================================================

\addto\captionsrussian{%
    \renewcommand{\chaptername}{Лекция}%
    \renewcommand{\contentsname}{Содержание}%
    \renewcommand{\figurename}{Рисунок}%
    \renewcommand{\tablename}{Таблица}%
    \renewcommand{\appendixname}{Приложение}%
    \renewcommand{\bibname}{Литература}%
}

\renewcommand{\proofname}{Доказательство}
\renewcommand{\qedsymbol}{\(\blacksquare\)}

% ============================================================
% ОФОРМЛЕНИЕ ЗАГОЛОВКОВ
% ============================================================

% Заголовок лекции
\titleformat{\chapter}[display]
    {\normalfont\bfseries}
    {\filleft
        \color{MaulRed}
        \fontsize{15}{18}\selectfont
        \MakeUppercase{\chaptername}\ \thechapter
    }
    {0.5cm}
    {\color{TextBlack}\Huge}
    [%
        \vspace{0.35cm}
        \color{MaulRed}
        \titlerule
    ]

\titleformat{name=\chapter,numberless}[display]
    {\normalfont\bfseries}
    {}
    {0pt}
    {\color{TextBlack}\Huge}
    [%
        \vspace{0.35cm}
        \color{MaulRed}
        \titlerule
    ]

\titlespacing*{\chapter}
    {0pt}
    {-15pt}
    {32pt}

% Разделы
\titleformat{\section}
    {\normalfont\Large\bfseries\color{TextBlack}}
    {\textcolor{MaulRed}{\thesection}}
    {0.75em}
    {}

\titlespacing*{\section}
    {0pt}
    {3.2ex plus 1ex minus 0.2ex}
    {1.3ex plus 0.2ex}

% Подразделы
\titleformat{\subsection}
    {\normalfont\large\bfseries\color{TextBlack}}
    {\textcolor{MaulRed}{\thesubsection}}
    {0.75em}
    {}

\titlespacing*{\subsection}
    {0pt}
    {2.6ex plus 0.8ex minus 0.2ex}
    {1ex plus 0.2ex}

% Пункты
\titleformat{\subsubsection}
    {\normalfont\normalsize\bfseries\color{DarkGray}}
    {\textcolor{MaulRed}{\thesubsubsection}}
    {0.7em}
    {}

% Абзацные заголовки
\titleformat{\paragraph}[runin]
    {\normalfont\normalsize\bfseries}
    {}
    {0pt}
    {}
    [.]

\titlespacing*{\paragraph}
    {0pt}
    {1.5ex}
    {0.6em}

% ============================================================
% КОЛОНТИТУЛЫ
% ============================================================

\setlength{\headheight}{26pt}

\pagestyle{fancy}
\fancyhf{}

\fancyhead[L]{%
    \small
    \nouppercase{\leftmark}
}

\fancyhead[R]{%
    \small
    \thepage
}

\renewcommand{\headrulewidth}{0.4pt}
\renewcommand{\footrulewidth}{0pt}

\renewcommand{\chaptermark}[1]{%
    \markboth{\chaptername\ \thechapter.\ #1}{}
}

\fancypagestyle{plain}{%
    \fancyhf{}
    \fancyfoot[C]{\thepage}
    \renewcommand{\headrulewidth}{0pt}
}

% ============================================================
% СПИСКИ
% ============================================================

\setlist[itemize,1]{
    label=\textcolor{MaulRed}{\textbullet},
    leftmargin=2em,
    itemsep=0.25em,
    topsep=0.5em
}

\setlist[itemize,2]{
    label=\textcolor{DarkGray}{\(\circ\)},
    leftmargin=2em,
    itemsep=0.15em,
    topsep=0.3em
}

\setlist[enumerate,1]{
    label=\textcolor{MaulRed}{\arabic*.},
    leftmargin=2.2em,
    itemsep=0.25em,
    topsep=0.5em
}

\setlist[enumerate,2]{
    label=\textcolor{DarkGray}{\alph*)},
    leftmargin=2em,
    itemsep=0.15em,
    topsep=0.3em
}

\setlist[description]{
    font=\normalfont\bfseries\color{MaulRed},
    style=nextline,
    leftmargin=2em,
    itemsep=0.4em
}

% ============================================================
% ПОДПИСИ К РИСУНКАМ И ТАБЛИЦАМ
% ============================================================

\captionsetup{
    font=small,
    labelfont={bf,color=MaulRed},
    labelsep=period,
    justification=centering,
    singlelinecheck=true
}

\captionsetup[subfigure]{
    font=small,
    labelfont=bf
}

% ============================================================
% НУМЕРАЦИЯ ФОРМУЛ И АЛГОРИТМОВ
% ============================================================

\numberwithin{equation}{chapter}
\numberwithin{algorithm}{chapter}

% ============================================================
% ТЕОРЕМНЫЕ СТИЛИ
% ============================================================

\newtheoremstyle{resultstyle}
    {9pt}
    {9pt}
    {\itshape}
    {}
    {\bfseries\color{MaulRed}}
    {.}
    {0.6em}
    {}

\newtheoremstyle{definitionstyle}
    {9pt}
    {9pt}
    {\normalfont}
    {}
    {\bfseries\color{MaulRed}}
    {.}
    {0.6em}
    {}

\newtheoremstyle{remarkstyle}
    {7pt}
    {7pt}
    {\normalfont}
    {}
    {\itshape\color{DarkGray}}
    {.}
    {0.6em}
    {}

% ------------------------------------------------------------
% Теоремы, леммы и следствия
% Используется единая сквозная нумерация внутри каждой лекции.
% ------------------------------------------------------------

\theoremstyle{resultstyle}

\newtheorem{theorem}{Теорема}[chapter]
\newtheorem{lemma}[theorem]{Лемма}
\newtheorem{proposition}[theorem]{Предложение}
\newtheorem{corollary}[theorem]{Следствие}
\newtheorem{claim}[theorem]{Утверждение}
\newtheorem{criterion}[theorem]{Критерий}

% ------------------------------------------------------------
% Определения, аксиомы и предположения
% ------------------------------------------------------------

\theoremstyle{definitionstyle}

\newtheorem{definition}[theorem]{Определение}
\newtheorem{axiom}[theorem]{Аксиома}
\newtheorem{assumption}[theorem]{Предположение}
\newtheorem{hypothesis}[theorem]{Гипотеза}
\newtheorem{condition}[theorem]{Условие}
\newtheorem{notation}[theorem]{Обозначение}
\newtheorem{convention}[theorem]{Соглашение}
\newtheorem{problem}[theorem]{Задача}
\newtheorem{example}[theorem]{Пример}
\newtheorem{counterexample}[theorem]{Контрпример}
\newtheorem{exercise}[theorem]{Упражнение}

% ------------------------------------------------------------
% Замечания и пояснения
% ------------------------------------------------------------

\theoremstyle{remarkstyle}

\newtheorem{remark}[theorem]{Замечание}
\newtheorem{intuition}[theorem]{Интуиция}
\newtheorem{warning}[theorem]{Предупреждение}
\newtheorem{observation}[theorem]{Наблюдение}
\newtheorem{discussion}[theorem]{Обсуждение}

% ------------------------------------------------------------
% Ненумерованные варианты
% ------------------------------------------------------------

\theoremstyle{resultstyle}

\newtheorem*{theorem*}{Теорема}
\newtheorem*{lemma*}{Лемма}
\newtheorem*{proposition*}{Предложение}
\newtheorem*{corollary*}{Следствие}

\theoremstyle{definitionstyle}

\newtheorem*{definition*}{Определение}
\newtheorem*{axiom*}{Аксиома}
\newtheorem*{assumption*}{Предположение}
\newtheorem*{example*}{Пример}

\theoremstyle{remarkstyle}

\newtheorem*{remark*}{Замечание}
\newtheorem*{intuition*}{Интуиция}
\newtheorem*{warning*}{Предупреждение}

% ============================================================
% ДОКАЗАТЕЛЬСТВА И РЕШЕНИЯ
% ============================================================

\newenvironment{proofidea}
    {\begin{proof}[Идея доказательства]}
    {\end{proof}}

\newenvironment{proofsketch}
    {\begin{proof}[Набросок доказательства]}
    {\end{proof}}

\newenvironment{solution}
    {\begin{proof}[Решение]}
    {\end{proof}}

\newenvironment{verification}
    {\begin{proof}[Проверка]}
    {\end{proof}}

% ============================================================
% СТРУКТУРА ОТДЕЛЬНОЙ ЛЕКЦИИ
% ============================================================

% Использование:
% \lecture{01}{Название лекции}
%
% или:
%
% \lecture{01}{Название лекции}[Дополнительная информация]

\NewDocumentCommand{\lecture}{m m O{}}{%
    \chapter{#2}%
    \label{lec:#1}%
    \ifstrempty{#3}{}{%
        \begin{flushright}
            \small\itshape\color{MiddleGray}
            #3
        \end{flushright}
        \vspace{0.5em}
    }%
}

% Цели лекции
\newenvironment{lecturegoals}
{%
    \par
    \medskip
    \noindent
    {\bfseries\color{MaulRed}Цели лекции}
    \begin{itemize}
}
{%
    \end{itemize}
    \medskip
}

% План лекции
\newenvironment{lectureplan}
{%
    \par
    \medskip
    \noindent
    {\bfseries\color{MaulRed}План лекции}
    \begin{enumerate}
}
{%
    \end{enumerate}
    \medskip
}

% Ключевые понятия
\newenvironment{keyconcepts}
{%
    \par
    \medskip
    \noindent
    {\bfseries\color{MaulRed}Ключевые понятия}
    \begin{itemize}
}
{%
    \end{itemize}
    \medskip
}

% Итоги лекции
\newenvironment{lecturesummary}
{%
    \begin{quote}
    \noindent
    {\bfseries\color{MaulRed}Итоги лекции.}\quad
}
{%
    \end{quote}
}

% Список обозначений конкретной лекции
\newenvironment{notationlist}
{%
    \par
    \medskip
    \noindent
    {\bfseries\color{MaulRed}Обозначения}
    \begin{description}
}
{%
    \end{description}
    \medskip
}

% Ссылка на исходную видеолекцию
\newcommand{\sourcevideo}[2]{%
    \par
    \medskip
    \noindent
    \textbf{\textcolor{MaulRed}{Видеолекция: }}
    \href{#1}{#2}
    \par
    \medskip
}

% Ссылка на дополнительный материал
\newcommand{\additionalsource}[2]{%
    \par
    \smallskip
    \noindent
    \textbf{Дополнительный источник: }
    \href{#1}{#2}
    \par
    \smallskip
}

% Заметка, которую нужно позднее доработать
\newcommand{\TODO}[1]{%
    \textcolor{MaulRed}{%
        \textbf{TODO: }#1%
    }%
}

% ============================================================
% АЛГОРИТМЫ И ПСЕВДОКОД НА РУССКОМ
% ============================================================

\floatname{algorithm}{Алгоритм}
\renewcommand{\listalgorithmname}{Список алгоритмов}

% ------------------------------------------------------------
% Базовые ключевые слова algpseudocode
% ------------------------------------------------------------

\algrenewcommand{\algorithmicrequire}{\textbf{Вход:}}
\algrenewcommand{\algorithmicensure}{\textbf{Выход:}}

\algrenewcommand{\algorithmicif}{\textbf{если}}
\algrenewcommand{\algorithmicthen}{\textbf{то}}
\algrenewcommand{\algorithmicelse}{\textbf{иначе}}

\algrenewcommand{\algorithmicfor}{\textbf{для}}
\algrenewcommand{\algorithmicforall}{\textbf{для всех}}
\algrenewcommand{\algorithmicdo}{\textbf{выполнять}}

\algrenewcommand{\algorithmicwhile}{\textbf{пока}}
\algrenewcommand{\algorithmicrepeat}{\textbf{повторять}}
\algrenewcommand{\algorithmicuntil}{\textbf{до выполнения условия}}

\algrenewcommand{\algorithmicreturn}{\textbf{вернуть}}
\algrenewcommand{\algorithmicprocedure}{\textbf{процедура}}
\algrenewcommand{\algorithmicfunction}{\textbf{функция}}
\algrenewcommand{\algorithmicend}{\textbf{конец}}

% Команда \ElsIf автоматически печатает:
% \algorithmicelse + \algorithmicif.
%
% Команды \EndIf, \EndFor и \EndWhile автоматически печатают:
% \algorithmicend + соответствующее ключевое слово.
%
% Поэтому отдельных команд algorithmicelsif,
% algorithmicendif, algorithmicendfor и
% algorithmicendwhile в algpseudocode нет.

\algrenewcommand{\algorithmiccomment}[1]{%
    \hfill
    \(\triangleright\)
    \textit{#1}%
}

% ------------------------------------------------------------
% Дополнительные команды для наших алгоритмов
% ------------------------------------------------------------

\algnewcommand{\Parameters}[1]{%
    \Statex
    \textbf{Параметры:} #1
}

\algnewcommand{\Initialize}[1]{%
    \Statex
    \textbf{Инициализация:} #1
}

\algnewcommand{\Output}[1]{%
    \Statex
    \textbf{Результат:} #1
}

\algnewcommand{\True}{%
    \textbf{истина}
}

\algnewcommand{\False}{%
    \textbf{ложь}
}

\algnewcommand{\Break}{%
    \State
    \textbf{выйти из цикла}
}

\algnewcommand{\Continue}{%
    \State
    \textbf{перейти к следующей итерации}
}

\algrenewcommand{\algorithmicindent}{1.5em}


% ============================================================
% ОФОРМЛЕНИЕ ПРОГРАММНОГО КОДА
% ============================================================

\lstdefinestyle{lecturecode}{
    basicstyle=\ttfamily\small,
    backgroundcolor=\color{CodeGray},
    frame=single,
    rulecolor=\color{MiddleGray},
    numbers=left,
    numberstyle=\scriptsize\color{MiddleGray},
    stepnumber=1,
    numbersep=8pt,
    showstringspaces=false,
    breaklines=true,
    breakatwhitespace=true,
    tabsize=4,
    columns=fullflexible,
    keepspaces=true,
    captionpos=b
}

\lstset{
    style=lecturecode
}

% ============================================================
% УМНЫЕ ССЫЛКИ
% ============================================================

\crefname{chapter}{лекция}{лекции}
\Crefname{chapter}{Лекция}{Лекции}

\crefname{section}{раздел}{разделы}
\Crefname{section}{Раздел}{Разделы}

\crefname{subsection}{подраздел}{подразделы}
\Crefname{subsection}{Подраздел}{Подразделы}

\crefname{equation}{формула}{формулы}
\Crefname{equation}{Формула}{Формулы}

\crefname{figure}{рисунок}{рисунки}
\Crefname{figure}{Рисунок}{Рисунки}

\crefname{table}{таблица}{таблицы}
\Crefname{table}{Таблица}{Таблицы}

\crefname{algorithm}{алгоритм}{алгоритмы}
\Crefname{algorithm}{Алгоритм}{Алгоритмы}

\crefname{theorem}{теорема}{теоремы}
\Crefname{theorem}{Теорема}{Теоремы}

\crefname{lemma}{лемма}{леммы}
\Crefname{lemma}{Лемма}{Леммы}

\crefname{proposition}{предложение}{предложения}
\Crefname{proposition}{Предложение}{Предложения}

\crefname{corollary}{следствие}{следствия}
\Crefname{corollary}{Следствие}{Следствия}

\crefname{claim}{утверждение}{утверждения}
\Crefname{claim}{Утверждение}{Утверждения}

\crefname{criterion}{критерий}{критерии}
\Crefname{criterion}{Критерий}{Критерии}

\crefname{definition}{определение}{определения}
\Crefname{definition}{Определение}{Определения}

\crefname{axiom}{аксиома}{аксиомы}
\Crefname{axiom}{Аксиома}{Аксиомы}

\crefname{assumption}{предположение}{предположения}
\Crefname{assumption}{Предположение}{Предположения}

\crefname{hypothesis}{гипотеза}{гипотезы}
\Crefname{hypothesis}{Гипотеза}{Гипотезы}

\crefname{condition}{условие}{условия}
\Crefname{condition}{Условие}{Условия}

\crefname{problem}{задача}{задачи}
\Crefname{problem}{Задача}{Задачи}

\crefname{example}{пример}{примеры}
\Crefname{example}{Пример}{Примеры}

\crefname{remark}{замечание}{замечания}
\Crefname{remark}{Замечание}{Замечания}

% ============================================================
% БАЗОВЫЕ ЧИСЛОВЫЕ МНОЖЕСТВА
% ============================================================

\newcommand{\NN}{\mathbb{N}}
\newcommand{\NNzero}{\mathbb{N}_0}
\newcommand{\ZZ}{\mathbb{Z}}
\newcommand{\QQ}{\mathbb{Q}}
\newcommand{\RR}{\mathbb{R}}
\newcommand{\CC}{\mathbb{C}}

\newcommand{\RRplus}{\mathbb{R}_{+}}
\newcommand{\RRplusplus}{\mathbb{R}_{++}}
\newcommand{\RRbar}{\overline{\mathbb{R}}}

% ============================================================
% СКОБКИ, НОРМЫ И МНОЖЕСТВА
% ============================================================

\DeclarePairedDelimiter{\abs}{\lvert}{\rvert}
\DeclarePairedDelimiter{\norm}{\lVert}{\rVert}
\DeclarePairedDelimiter{\ceil}{\lceil}{\rceil}
\DeclarePairedDelimiter{\floor}{\lfloor}{\rfloor}
\DeclarePairedDelimiter{\paren}{(}{)}
\DeclarePairedDelimiter{\bracks}{[}{]}
\DeclarePairedDelimiter{\angles}{\langle}{\rangle}

\newcommand{\set}[1]{%
    \left\{#1\right\}
}

\newcommand{\setgiven}[2]{%
    \left\{
        #1
        \,\middle|\,
        #2
    \right\}
}

\newcommand{\suchthat}{%
    \;\middle|\;
}

\newcommand{\card}[1]{%
    \abs*{#1}
}

\newcommand{\inner}[2]{%
    \left\langle
        #1,#2
    \right\rangle
}

\newcommand{\dualpair}[2]{%
    \left\langle
        #1,#2
    \right\rangle
}

\newcommand{\weightednorm}[2]{%
    \norm{#1}_{#2}
}

\newcommand{\dualnorm}[1]{%
    \norm{#1}_{*}
}

\newcommand{\seminorm}[1]{%
    \left|#1\right|
}

\newcommand{\evalat}[2]{%
    \left.
        #1
    \right|_{#2}
}

\newcommand{\restr}[2]{%
    #1\!\restriction_{#2}
}

% ============================================================
% ВЕКТОРЫ И МАТРИЦЫ
% ============================================================

\newcommand{\vct}[1]{\bm{#1}}
\newcommand{\mat}[1]{\bm{#1}}

\newcommand{\zero}{\mathbf{0}}
\newcommand{\one}{\mathbf{1}}
\newcommand{\Id}{\mathbf{I}}

\newcommand{\trans}{^{\mathsf{T}}}
\newcommand{\herm}{^{\mathsf{H}}}
\newcommand{\inv}{^{-1}}
\newcommand{\pinv}{^{\dagger}}

\newcommand{\kron}{\otimes}

\newcommand{\vectwo}[2]{%
    \begin{pmatrix}
        #1\\
        #2
    \end{pmatrix}
}

\newcommand{\vecthree}[3]{%
    \begin{pmatrix}
        #1\\
        #2\\
        #3
    \end{pmatrix}
}

% ============================================================
% ЛИНЕЙНАЯ АЛГЕБРА
% ============================================================

\DeclareMathOperator{\trace}{tr}
\DeclareMathOperator{\rank}{rank}
\DeclareMathOperator{\diag}{diag}
\DeclareMathOperator{\blkdiag}{blkdiag}
\DeclareMathOperator{\Ker}{Ker}
\DeclareMathOperator{\Image}{Im}
\DeclareMathOperator{\Range}{range}
\DeclareMathOperator{\Null}{null}
\DeclareMathOperator{\Lin}{Lin}
\DeclareMathOperator{\End}{End}
\DeclareMathOperator{\linspan}{span}
\DeclareMathOperator{\spec}{spec}
\DeclareMathOperator{\spectralradius}{\rho}
\DeclareMathOperator{\cond}{cond}
\DeclareMathOperator{\vecop}{vec}

\newcommand{\lambdaMin}[1]{%
    \lambda_{\min}\!\left(#1\right)
}

\newcommand{\lambdaMax}[1]{%
    \lambda_{\max}\!\left(#1\right)
}

\newcommand{\sigmaMin}[1]{%
    \sigma_{\min}\!\left(#1\right)
}

\newcommand{\sigmaMax}[1]{%
    \sigma_{\max}\!\left(#1\right)
}

\newcommand{\conditionnumber}[1]{%
    \kappa\!\left(#1\right)
}

\newcommand{\Sym}[1]{%
    \mathbb{S}^{#1}
}

\newcommand{\PSD}[1]{%
    \mathbb{S}^{#1}_{+}
}

\newcommand{\PD}[1]{%
    \mathbb{S}^{#1}_{++}
}

% ============================================================
% ПРОИЗВОДНЫЕ И ДИФФЕРЕНЦИАЛЫ
% ============================================================

\newcommand{\dd}{%
    \mathop{}\!\mathrm{d}
}

\newcommand{\od}[2]{%
    \frac{\dd #1}{\dd #2}
}

\newcommand{\odn}[3]{%
    \frac{\dd^{#3} #1}{\dd #2^{#3}}
}

\newcommand{\pd}[2]{%
    \frac{\partial #1}{\partial #2}
}

\newcommand{\pdn}[3]{%
    \frac{\partial^{#3} #1}{\partial #2^{#3}}
}

\newcommand{\grad}{\nabla}
\newcommand{\hess}{\nabla^{2}}
\newcommand{\jacobian}{\mathrm{D}}

\newcommand{\gradat}[2]{%
    \nabla #1\!\left(#2\right)
}

\newcommand{\hessat}[2]{%
    \nabla^{2} #1\!\left(#2\right)
}

% ============================================================
% ПРЕДЕЛЫ И АСИМПТОТИКА
% ============================================================

\newcommand{\BigO}{\mathcal{O}}
\newcommand{\littleo}{o}
\newcommand{\BigTheta}{\Theta}
\newcommand{\BigOmega}{\Omega}
\newcommand{\softO}{\widetilde{\mathcal{O}}}

\newcommand{\asympeq}{\sim}

% ============================================================
% БАЗОВЫЕ МАТЕМАТИЧЕСКИЕ ОПЕРАТОРЫ
% ============================================================

\DeclareMathOperator*{\argmin}{arg\,min}
\DeclareMathOperator*{\argmax}{arg\,max}

\DeclareMathOperator*{\esssup}{ess\,sup}
\DeclareMathOperator*{\essinf}{ess\,inf}

\DeclareMathOperator{\sign}{sign}
\DeclareMathOperator{\supp}{supp}
\DeclareMathOperator{\diam}{diam}
\DeclareMathOperator{\dist}{dist}

\DeclareMathOperator{\interior}{int}
\DeclareMathOperator{\closure}{cl}
\DeclareMathOperator{\boundary}{bd}
\DeclareMathOperator{\relint}{ri}

\DeclareMathOperator{\dom}{dom}
\DeclareMathOperator{\epi}{epi}
\DeclareMathOperator{\hypo}{hypo}
\DeclareMathOperator{\graphop}{gph}

\DeclareMathOperator{\conv}{conv}
\DeclareMathOperator{\cone}{cone}
\DeclareMathOperator{\aff}{aff}

\DeclareMathOperator{\Fix}{Fix}
\DeclareMathOperator{\zer}{zer}

\DeclareMathOperator{\proxop}{prox}
\DeclareMathOperator{\projop}{proj}

\DeclareMathOperator{\softmax}{softmax}
\DeclareMathOperator{\ReLU}{ReLU}

% ============================================================
% ЭКСПОНЕНТА, МНИМАЯ ЕДИНИЦА И ОПРЕДЕЛЯЮЩЕЕ РАВЕНСТВО
% ============================================================

\newcommand{\e}{\mathrm{e}}
\newcommand{\ii}{\mathrm{i}}

\newcommand{\defeq}{\coloneqq}
\newcommand{\eqdef}{\eqqcolon}

% ============================================================
% ОПТИМИЗАЦИОННЫЕ ЗАДАЧИ
% ============================================================

\newcommand{\minimize}{%
    \operatorname*{minimize}
}

\newcommand{\maximize}{%
    \operatorname*{maximize}
}

\newcommand{\subjectto}{%
    \text{при условиях}
}

\newcommand{\find}{%
    \operatorname*{find}
}

\newcommand{\optval}{%
    f^{\star}
}

\newcommand{\xopt}{%
    x^{\star}
}

\newcommand{\Xopt}{%
    X^{\star}
}

\newcommand{\Lagrangian}{%
    \mathcal{L}
}

\newcommand{\KKT}{%
    \mathrm{KKT}
}

% ============================================================
% ВЫПУКЛЫЙ АНАЛИЗ
% ============================================================

\newcommand{\subdiff}[2][]{%
    \partial_{#1} #2
}

\newcommand{\prox}[2]{%
    \operatorname{prox}_{#1}\!\left(#2\right)
}

\newcommand{\projection}[2]{%
    \operatorname{proj}_{#1}\!\left(#2\right)
}

\newcommand{\distto}[2]{%
    \operatorname{dist}\!\left(#1,#2\right)
}

\newcommand{\normalcone}[2]{%
    N_{#1}\!\left(#2\right)
}

\newcommand{\tangentcone}[2]{%
    T_{#1}\!\left(#2\right)
}

\newcommand{\indicatorfunction}[1]{%
    \delta_{#1}
}

\newcommand{\indevent}[1]{%
    \mathbf{1}_{\set{#1}}
}

\newcommand{\fenchelconjugate}[1]{%
    #1^{*}
}

\newcommand{\biconjugate}[1]{%
    #1^{**}
}

\newcommand{\bregman}[3]{%
    D_{#1}\!\left(#2,#3\right)
}

\newcommand{\simplex}[1]{%
    \Delta^{#1}
}

\newcommand{\ball}[2]{%
    \mathbb{B}_{#1}\!\left(#2\right)
}

\newcommand{\unitball}{%
    \mathbb{B}
}

\newcommand{\positivepart}[1]{%
    \left[#1\right]_{+}
}

\newcommand{\negativepart}[1]{%
    \left[#1\right]_{-}
}

% ============================================================
% ВЕРОЯТНОСТЬ И СТАТИСТИКА
% ============================================================

\providecommand{\Prob}{}
\renewcommand{\Prob}{\mathbb{P}}
\providecommand{\Expect}{}
\renewcommand{\Expect}{\mathbb{E}}

\NewDocumentCommand{\E}{o m}{%
    \mathbb{E}%
    \IfValueT{#1}{_{#1}}%
    \!\left[#2\right]
}

\NewDocumentCommand{\Prb}{o m}{%
    \mathbb{P}%
    \IfValueT{#1}{_{#1}}%
    \!\left(#2\right)
}

\DeclareMathOperator{\Var}{Var}
\DeclareMathOperator{\Cov}{Cov}
\DeclareMathOperator{\Corr}{Corr}
\DeclareMathOperator{\Law}{Law}

\newcommand{\given}{%
    \,\middle|\,
}

\newcommand{\VarOf}[1]{%
    \operatorname{Var}\!\left(#1\right)
}

\newcommand{\CovOf}[2]{%
    \operatorname{Cov}\!\left(#1,#2\right)
}

\newcommand{\NormalDist}{\mathcal{N}}
\newcommand{\UniformDist}{\mathcal{U}}
\newcommand{\BernoulliDist}{\operatorname{Bern}}
\newcommand{\CategoricalDist}{\operatorname{Cat}}

\newcommand{\iid}{%
    \mathrel{\overset{\mathrm{i.i.d.}}{\sim}}
}

\newcommand{\independent}{%
    \mathrel{\perp\!\!\!\perp}
}

\newcommand{\almostsure}{%
    \text{п.н.}
}

\newcommand{\asconv}{%
    \xrightarrow{\text{п.н.}}
}

\newcommand{\pconv}{%
    \xrightarrow{\Prob}
}

\newcommand{\dconv}{%
    \xrightarrow{d}
}

\newcommand{\loneconv}{%
    \xrightarrow{L^{1}}
}

\newcommand{\ltwoconv}{%
    \xrightarrow{L^{2}}
}

% ============================================================
% ИНФОРМАЦИОННЫЕ ВЕЛИЧИНЫ
% ============================================================

\newcommand{\Entropy}[1]{%
    H\!\left(#1\right)
}

\newcommand{\ConditionalEntropy}[2]{%
    H\!\left(#1\mid #2\right)
}

\newcommand{\MutualInformation}[2]{%
    I\!\left(#1;#2\right)
}

\newcommand{\KL}[2]{%
    D_{\mathrm{KL}}\!\left(
        #1\,\middle\|\,#2
    \right)
}

% ============================================================
% МАШИННОЕ ОБУЧЕНИЕ
% ============================================================

\newcommand{\Dataset}{\mathcal{D}}
\newcommand{\Samples}{\mathcal{S}}
\newcommand{\HypothesisClass}{\mathcal{H}}

\newcommand{\Loss}{\ell}
\newcommand{\Risk}{\mathcal{R}}
\newcommand{\EmpiricalRisk}{\widehat{\mathcal{R}}}

\newcommand{\parameters}{\theta}
\newcommand{\parameterSpace}{\Theta}

\newcommand{\prediction}{\widehat{y}}
\newcommand{\features}{x}
\newcommand{\labelvalue}{y}

\newcommand{\model}{f_{\theta}}

\newcommand{\regularizer}{r}

\newcommand{\batch}{\mathcal{B}}
\newcommand{\batchsize}{B}

\newcommand{\learningrate}{\eta}

\newcommand{\stochasticgradient}{%
    g
}

\newcommand{\gradientestimate}{%
    \widehat{\nabla f}
}

% ============================================================
% РАСПРЕДЕЛЁННАЯ ОПТИМИЗАЦИЯ
% ============================================================

\newcommand{\Graph}{\mathcal{G}}
\newcommand{\Vertices}{\mathcal{V}}
\newcommand{\Edges}{\mathcal{E}}

\newcommand{\numagents}{n}

\newcommand{\neighbors}[1]{%
    \mathcal{N}_{#1}
}

\newcommand{\degreeof}[1]{%
    d_{#1}
}

\newcommand{\Adjacency}{\mathbf{A}}
\newcommand{\DegreeMatrix}{\mathbf{D}}
\newcommand{\Laplacian}{\mathbf{L}}
\newcommand{\MixingMatrix}{\mathbf{W}}

\newcommand{\IncidenceMatrix}{\mathbf{B}}

\newcommand{\localobjective}[1]{%
    f_{#1}
}

\newcommand{\globalobjective}{%
    F
}

\newcommand{\localvariable}[1]{%
    x_{#1}
}

\newcommand{\stackedvariable}{%
    \mathbf{x}
}

\newcommand{\networkaverage}[1]{%
    \overline{#1}
}

\newcommand{\ConsensusProjector}[1]{%
    \frac{1}{#1}\one\one\trans
}

\newcommand{\DisagreementProjector}[1]{%
    \Id-\frac{1}{#1}\one\one\trans
}

\newcommand{\ConsensusSpace}{%
    \linspan\set{\one}
}

\newcommand{\consensuserror}[1]{%
    \norm{
        #1-
        \ConsensusProjector{\numagents}#1
    }
}

\newcommand{\spectralgap}[1]{%
    1-\lambda_{2}\!\left(#1\right)
}

\newcommand{\networkcondition}[1]{%
    \frac{
        \lambda_{\max}\!\left(#1\right)
    }{
        \lambda_{\min}^{+}\!\left(#1\right)
    }
}

% ============================================================
% ПОСЛЕДОВАТЕЛЬНОСТИ И ИТЕРАЦИИ
% ============================================================

\newcommand{\sequence}[2]{%
    \set{#1}_{#2}
}

\newcommand{\iter}[2]{%
    #1^{#2}
}

\newcommand{\nextiter}[1]{%
    #1^{k+1}
}

\newcommand{\currentiter}[1]{%
    #1^{k}
}

\newcommand{\previousiter}[1]{%
    #1^{k-1}
}

\newcommand{\average}[1]{%
    \overline{#1}
}

\newcommand{\runningaverage}[2]{%
    \frac{1}{#1}
    \sum_{k=1}^{#1}
    #2
}

% ============================================================
% ГЛАДКОСТЬ, ВЫПУКЛОСТЬ И ЧИСЛО ОБУСЛОВЛЕННОСТИ
% ============================================================

\newcommand{\smoothnessconstant}{L}
\newcommand{\strongconvexityconstant}{\mu}

\newcommand{\conditionratio}{%
    \kappa
    \defeq
    \frac{L}{\mu}
}

\newcommand{\LipschitzConstant}{L}
\newcommand{\StrongConvexityConstant}{\mu}

% ============================================================
% ТИПОВЫЕ ОПТИМИЗАЦИОННЫЕ ОТОБРАЖЕНИЯ
% ============================================================

\newcommand{\gradientstep}[2]{%
    #1-\learningrate\nabla #2
}

\newcommand{\resolvent}[2]{%
    \left(
        \Id+#1#2
    \right)^{-1}
}

\newcommand{\fixedpointset}[1]{%
    \operatorname{Fix}\!\left(#1\right)
}

\newcommand{\zeroset}[1]{%
    \operatorname{zer}\!\left(#1\right)
}

% ============================================================
% СПЕЦИАЛЬНЫЕ ОБОЗНАЧЕНИЯ
% ============================================================

\newcommand{\eps}{\varepsilon}

\newcommand{\weakto}{%
    \rightharpoonup
}

\newcommand{\weakstarto}{%
    \overset{*}{\rightharpoonup}
}

\newcommand{\compose}{\circ}

\newcommand{\directsum}{\oplus}

\newcommand{\bigdirectsum}{\bigoplus}

% ============================================================
% ВЫДЕЛЕНИЕ ВАЖНЫХ ФОРМУЛ
% ============================================================

\newcommand{\importantformula}[1]{%
    \boxed{\displaystyle #1}
}

% ============================================================
% TIKZ: СТИЛИ ДЛЯ ГРАФОВ И АЛГОРИТМОВ
% ============================================================

\tikzset{
    vertex/.style={
        circle,
        draw=MaulRed,
        thick,
        minimum size=8mm,
        inner sep=1pt,
        fill=white
    },
    active vertex/.style={
        vertex,
        fill=MaulRed,
        text=white
    },
    edge/.style={
        thick,
        draw=DarkGray
    },
    directed edge/.style={
        edge,
        -{Latex[length=2mm]}
    },
    communication edge/.style={
        thick,
        dashed,
        draw=MiddleGray
    },
    block/.style={
        rectangle,
        draw=MaulRed,
        thick,
        rounded corners=2pt,
        minimum width=2.8cm,
        minimum height=0.9cm,
        align=center
    },
    decision/.style={
        diamond,
        draw=MaulRed,
        thick,
        aspect=2,
        align=center
    },
    flow arrow/.style={
        thick,
        -{Latex[length=2mm]}
    }
}

% ============================================================
% КОНЕЦ PREAMBLE.TEX
% ============================================================